% ===================================================================
% 05-frontend.tex — Frontend SvelteKit
% Źródła: frontend/README.md, docs/frontend/{routes,library,components}.md
% oraz kod frontend/src (regiony, urls, clientErrors, healthz, backend.ts).
% ===================================================================

\chapter{Frontend: struktura i trasy}\label{ch:fe-trasy}

Frontend PlayLink to aplikacja SvelteKit realizująca logowanie fraza BIP39 oraz przeglądarkę
pokoi (lobby). Komunikuje się z backendem FastAPI po REST i WebSocket. Pełny opis kontraktu API
i protokołów czasu rzeczywistego znajduje się w rozdziałach o backendzie; tutaj opisano
wyłącznie stronę klienta.

\section{Stos technologiczny i uruchomienie}\label{sec:fe-stos}

\begin{itemize}[nosep]
  \item \textbf{Svelte} \code{^5.51.0} w trybie runes (\code{$state}, \code{$derived}, \code{$effect}, \code{$props}),
  \item \textbf{SvelteKit} \code{^2.50.2} z \code{@sveltejs/adapter-node},
  \item \textbf{TypeScript} (tryb strict),
  \item \textbf{Bun} --- środowisko uruchomieniowe, menedżer pakietów i bundler (npm, pnpm i yarn nie są używane),
  \item \textbf{ethers} + \textbf{viem} --- lokalna derywacja z mnemoniku i podpisy ECDSA,
  \item \textbf{jwt-decode} --- dekodowanie ciasteczka sesji po stronie serwera.
\end{itemize}

Polecenia: \code{bun install}, \code{bun dev} (serwer deweloperski, domyślnie
\code{http://localhost:5173/}), \code{bun run build} (wynik: serwer Node dla dołączonego
\code{Dockerfile}), \code{bun run preview}. Kontrola jakości: \code{bun run check}
(svelte-check + TypeScript), \code{bun run lint} (ESLint), \code{bun run format} (Prettier).
Zalecane środowisko: VS Code z rozszerzeniem Svelte.

\subsection{Zmienne środowiskowe}

Plik \code{frontend/.env} powstaje z \code{frontend/.env.example}. Zmienne publiczne są
odczytywane przez \code{$env/dynamic/public}, więc wartości można zmieniać w czasie działania
kontenera, bez przebudowy (wdrożenie wieloregionowe --- patrz część o operacjach).

\begin{longtable}{@{}p{0.34\textwidth}p{0.62\textwidth}@{}}
\toprule
\textbf{Zmienna} & \textbf{Znaczenie} \\
\midrule
\endhead
\code{PUBLIC_BACKEND_URL} & Bazowy URL backendu HTTP (np. \code{http://localhost:8000}); używana w przeglądarce. \\
\code{PUBLIC_WS_URL} & Bazowy URL WebSocket (\code{ws://} lub \code{wss://}), np. \code{ws://localhost:8000}. \\
\code{BACKEND_INTERNAL_URL} & (prywatna) adres backendu widziany z kontenera frontendu, np. nazwa usługi w compose; używany przez loadery i form actions po stronie serwera. \\
\code{PUBLIC_REGION} & Region obsługiwany przez daną instancję frontendu: \code{eu} lub \code{us} (każda inna wartość oznacza \code{eu}). \\
\code{PUBLIC_ENTRY_HOST} & Host (opcjonalnie z portem), który zamiast aplikacji renderuje wybór regionu. Dla lokalnej uprzęży dwuregionowej \code{localhost:3000}; domyślnie w \code{.env.example} \code{localhost:9999}, aby zwykły serwer deweloperski nie ukrywał aplikacji za wyborem. \\
\code{PUBLIC_REGION_EU_FRONTEND}, \code{PUBLIC_REGION_EU_API}, \code{PUBLIC_REGION_EU_WS} & Publiczne adresy frontendu, API i WebSocket regionu EU (lokalnie: \code{http://localhost:3000}, \code{http://localhost:8000}, \code{ws://localhost:8000}). \\
\code{PUBLIC_REGION_US_FRONTEND}, \code{PUBLIC_REGION_US_API}, \code{PUBLIC_REGION_US_WS} & Odpowiedniki dla regionu US (lokalnie: porty 3001 i 8001). \\
\bottomrule
\end{longtable}

W \code{.env.example} zaznaczono też, że adapter-node ma używać oryginalnego nagłówka
\code{Host} i schematu z \code{X-Forwarded-Proto}; zmiennej \code{ORIGIN} nie ustawia się,
gdy jeden kontener frontendu obsługuje zarówno host wejściowy, jak i hosty regionalne
(ustawienie w \code{compose.yaml}).

\section{Struktura katalogów}\label{sec:fe-struktura}

\begin{lstlisting}
frontend/src/
  app.html, app.d.ts          # szablon HTML; typ App.Locals
  hooks.server.ts             # handle (sesja) + handleError (log)
  hooks.client.ts             # raportowanie bledow przegladarki
  lib/
    auth.ts, signing.ts, signingKey.ts       # tozsamosc i podpisy
    roomsStore.ts, chatStore.ts              # store'y WebSocket
    hintsContext.svelte.ts                   # kontekst podpowiedzi klawiszy
    lobbyLocations.ts                        # lokalizacje lobby, Haversine
    regions.ts, urls.ts, clientErrors.ts     # regiony, URL-e, bledy klienta
    components/                              # komponenty domenowe
      chrome/                                # UI kit w stylu Diablo II (18)
    server/backend.ts, server/log.ts         # fetch SSR, log JSON
    styles/tokens.css, styles/chrome.css, global.css
  routes/
    +layout.svelte, +layout.server.ts
    +page.svelte  about/  auth/  profile/  rooms/  rooms/[name]/  test/
    api/client-errors/+server.ts   healthz/+server.ts
  test/                        # testy jednostkowe (Vitest)
\end{lstlisting}

\section{Model BFF i sesja}\label{sec:fe-bff}

Frontend stosuje wzorzec Backend-for-Frontend. Token JWT wydany przez backend nigdy nie trafia
do pamięci JavaScriptu w przeglądarce w trwały sposób: klient przesyła go jednorazowo do form
action \code{?/login}, a serwer zapisuje wyłącznie w ciasteczku \code{session} o atrybutach
\code{httpOnly}, \code{path=/}, \code{maxAge} 3600 s (1 godzina). Strony otrzymują tylko dane
pochodne (adres, nazwa użytkownika, \code{isAdmin}).

\begin{enumerate}[nosep]
  \item Klient generuje albo wprowadza frazę BIP39.
  \item \code{authenticate(mnemonic)} z \code{$lib/auth.ts} wykonuje trójfazowy handshake z backendem
        (żądanie nonce, lokalny podpis, weryfikacja; szczegóły w rozdziale o bibliotece klienta).
  \item Backend zwraca \code{{ token, username }}.
  \item Klient wysyła \code{token} jako pole formularza do akcji \code{?/login}.
  \item Akcja ustawia ciasteczko \code{session}; atrybut \code{secure} zależy od protokołu
        żądania (\code{https:}).
  \item Wszystkie loadery i akcje serwerowe czytają \code{cookies.get('session')} i dekodują je \code{jwtDecode}.
  \item \code{hooks.server.ts} przy każdym żądaniu dekoduje JWT do \code{event.locals.user}
        (\code{{ address, isAdmin }}); wygasłe lub nieprawidłowe ciasteczko jest usuwane.
  \item Wylogowanie: akcja \code{?/logout} wykonuje \code{cookies.delete('session', { path: '/' })}.
\end{enumerate}

Klucz podpisujący (klucz prywatny z mnemoniku) jest przechowywany w \code{sessionStorage}
(\code{saveSigningKey()}), aby podpisywać wiadomości czatu bez ponownego wpisywania frazy;
znika po wylogowaniu lub zamknięciu karty.

Ciasteczko sesji jest \emph{host-only}, więc nie jest współdzielone między frontendami
regionalnymi: przełączenie regionu wymaga ponownego logowania tą samą frazą 12 słów
(patrz rozdział o regionach w UI).

\section{Mapa tras}\label{sec:fe-mapa}

\begin{longtable}{@{}p{0.17\textwidth}p{0.45\textwidth}p{0.31\textwidth}@{}}
\toprule
\textbf{Trasa} & \textbf{Przeznaczenie} & \textbf{Dostęp} \\
\midrule
\endhead
\code{/} & Ekran powitalny: herb, tytuł, przycisk „Enter Realm” & Brak ograniczeń \\
\code{/about} & Statyczna strona informacyjna & Brak ograniczeń \\
\code{/auth} & Logowanie fraza BIP39 (3 fazy), akcje login/logout & Brak; wygląd zależy od sesji \\
\code{/profile} & Podgląd i zmiana nazwy użytkownika, sigil, adres, daty & Serwer przekierowuje na \code{/auth} bez sesji \\
\code{/rooms} & Lista pokoi w czasie rzeczywistym, filtry, tworzenie/dołączanie/opuszczanie, administracja grami & Load publiczny (gry, lokalizacje); akcje wymagają sesji \\
\code{/rooms/[name]} & Szczegóły pokoju: czat (\code{ChatStore}), lista członków, wydarzenia i RSVP, wyrzucanie/zamykanie & Wymagana sesja i członkostwo (lub administrator) \\
\code{/test} & Obserwatorium deweloperskie: status backendu, opóźnienie, surowy JSON & Brak ograniczeń \\
\code{/healthz} & Sonda żywotności frontendu (\code{{"status":"ok"}}) & Publiczna \\
\code{/api/client-errors} & Odbiór błędów przeglądarki (POST) & Publiczna, z limitem \\
\bottomrule
\end{longtable}

Trasy \code{/healthz} i \code{/api/client-errors} to endpointy serwerowe, nie strony; opisano je
w rozdziale o bibliotece klienta.

\section{Układ główny (layout)}\label{sec:fe-layout}

\code{routes/+layout.server.ts} zwraca dla każdej strony: \code{isAuthenticated}
(obecność ciasteczka \code{session}), \code{isEntryHost} (czy host żądania równa się
\code{PUBLIC_ENTRY_HOST}), \code{region} (\code{PUBLIC_REGION}) oraz listę \code{regions}.
\code{+layout.svelte} renderuje wtedy jeden z dwóch widoków:

\begin{itemize}[nosep]
  \item \textbf{Host wejściowy} (\code{isEntryHost}): w ramce \code{PageFrame} wyłącznie komponent
        \code{RegionChooser} --- wybór regionu.
  \item \textbf{Host regionalny}: pasek marki z odznaką aktywnego regionu (np. \code{EU}),
        link „$\rightarrow$ druga strona” do tej samej ścieżki w drugim regionie
        (z tytułem: osobna sesja, logowanie tą samą frazą), zakładki
        (Identity, Rooms, Profile --- tylko po zalogowaniu, About), baner
        \code{RegionProposalBanner}, treść strony oraz stopka \code{HintBar}
        zasilana kontekstem podpowiedzi (domyślnie \code{Esc: Back}, \code{?: Help}).
\end{itemize}

\section{Trasy szczegółowo}\label{sec:fe-trasy-szczegoly}

\subsection{\texttt{/} --- ekran powitalny}

Strona w pełni statyczna (bez \code{+page.server.ts}): środkowy herb \code{Crest}
(\code{size=104}, \code{tone="gold"}), tytuł „PLAYLINK”, podtytuł „Lobbies · Rooms · Kin”
i przycisk \code{OrnateButton} prowadzący do \code{/auth}. Klawisz \code{Enter} także przechodzi
do \code{/auth}; \code{HintBar} pokazuje \code{Enter: Enter Realm}.

\subsection{\texttt{/about}}

Statyczna strona wyjaśniająca model tożsamości (niekustodialny, oparty o BIP39) oraz „rytuał”
w trzech krokach: (1) \emph{Forge a Lobby} --- wybór gry, liczby graczy i nazwy;
(2) \emph{Watch the Hall} --- otwarte lobby spływają na listę w czasie rzeczywistym, z filtrami
slotów, statusu i pingu; (3) \emph{Stand in the Room} --- dołączenie, czat, koordynacja,
opuszczenie lub wygaśnięcie pokoju po czasie. \code{HintBar}: \code{Esc: Back}.

\subsection{\texttt{/auth} --- tożsamość}

\paragraph{Load} (\code{auth/+page.server.ts}) zwraca \code{{ user: { address, username } | null }}.
Brak ciasteczka daje \code{null}; ciasteczko jest dekodowane (\code{sub}, \code{username},
\code{exp}); nieprawidłowe lub wygasłe jest usuwane.

\paragraph{Form actions.}

\begin{longtable}{@{}p{0.12\textwidth}p{0.22\textwidth}p{0.55\textwidth}@{}}
\toprule
\textbf{Akcja} & \textbf{Wejście} & \textbf{Działanie} \\
\midrule
\endhead
\code{login} & \code{token} (string) & Bez wywołania backendu (BFF tylko na ciasteczku). Ustawia ciasteczko \code{session}: \code{httpOnly}, \code{sameSite=lax}, \code{secure} dla \code{https:}, \code{maxAge} 3600 s, \code{path=/}. Pusty token daje \code{{ success: false, error }}; sukces: \code{{ success: true }}. \\
\code{logout} & --- & Usuwa ciasteczko \code{session}; zwraca \code{{ success: true }}. \\
\bottomrule
\end{longtable}

Wybór \code{lax} (zamiast \code{strict}) jest celowy: pozwala działać linkom
przychodzącym z zewnątrz (np. z Discorda); ciasteczko jest host-only i nie jest współdzielone
z drugim frontendem regionalnym.

\paragraph{Interfejs} zależy od \code{data.user}.

\emph{Niezalogowany (Vault Keeper).} Siatka \code{MnemonicInput} (3 kolumny) na 12 słów BIP39.
Skróty: \code{G} --- wygeneruj nową frazę (\code{generateMnemonic()}), \code{C} --- kopiuj frazę
(awaryjnie \code{document.execCommand('copy')}, gdy asynchroniczny schowek jest zablokowany
w kontekście HTTP), \code{Enter} --- zatwierdź. Funkcja \code{startAuth()}:
\begin{enumerate}[nosep]
  \item woła \code{authenticate(mnemonic)} (faza 1: \code{POST /auth/request-nonce?address=...},
        limit żądań: 429 daje komunikat „Rate limited”; faza 2: derywacja \code{HDNodeWallet},
        zapis klucza w \code{sessionStorage}, podpis wiadomości
        \code{"Sign in to Playlink\nNonce: <nonce>"}; faza 3: \code{POST /auth/verify}
        z \code{{ address, nonce, signature }}),
  \item wysyła JWT do \code{?/login},
  \item po sukcesie wywołuje \code{invalidateAll()}, by przeładować aktywne loadery.
\end{enumerate}
Błędy są pokazywane w \code{SystemDialog}.

\emph{Zalogowany (Active Covenant).} Sigil (88 px), nazwa użytkownika i adres. Przycisk
„Sever Covenant” (klawisz \code{L}) wysyła \code{?/logout}; po wylogowaniu \code{invalidateAll()}
przełącza stronę do widoku niezalogowanego.

\subsection{\texttt{/profile} --- profil}

\paragraph{Load.} Brak ciasteczka: \code{redirect(303, '/auth')}. W przeciwnym razie
\code{GET /users/me} z nagłówkiem \code{Authorization: Bearer <session>} (przez
\code{backendFetch}); odpowiedź inna niż OK także przekierowuje na \code{/auth}. Zwraca
\code{{ identity_address, username, created_at, last_login }}.

\paragraph{Akcja \texttt{update}.} Wejście: \code{username} (przycinany, 3--20 znaków,
\code{[a-zA-Z0-9_-]}); wywołanie \code{PATCH /users/me} z JSON \code{{ username }}.
Sukces: \code{{ success: true, username }}. Błędy: brak sesji \code{fail(401, 'Not authenticated')},
pusta nazwa \code{fail(400, 'Username is required.')}, błąd serwera \code{fail(500, 'Server error')},
pozostałe kody przekazywane z backendu.

\paragraph{Interfejs.} Sigil (88 px), nazwa, \code{identity_address}, \code{created_at}
i \code{last_login} (\code{toLocaleString()}, „---” dla \code{null}). Pole nazwy z walidacją
kliencką \code{^[a-zA-Z0-9_-]{3,20}$}; przycisk zapisu jest wyłączony, gdy
\code{!dirty || !clientValid || saving}. Formularz używa \code{use:enhance}; błędy serwera
w \code{SystemDialog}. \code{HintBar}: \code{Enter: Save Name}.

\subsection{\texttt{/rooms} --- lista gier (pokoi)}

\paragraph{Load} (\code{rooms/+page.server.ts}) zwraca \code{isAuthenticated}, \code{user},
\code{isAdmin}, \code{games} (\code{string[]}), \code{lobbyLocations} i \code{defaultLobbyLocation}.
Ciasteczko jest dekodowane po cichu (błąd nie przerywa). Równolegle pobierane są
\code{GET /games} (przy błędzie \code{[]}) i \code{GET /lobby-locations}
(odpowiedź \code{{ default?, locations?: { code, label, lat, lon }[] }} jest parsowana
przez \code{parseLobbyLocations()} z awaryjnym \code{FALLBACK_LOBBY_LOCATIONS}).
Load nigdy nie przekierowuje --- strona renderuje się również bez sesji.

\paragraph{Form actions} (wszystkie wymagają sesji, w przeciwnym razie \code{fail(401)};
sukces zwraca \code{{ success: true, message }}, błąd \code{fail(status, { error })}):

\begin{longtable}{@{}p{0.14\textwidth}p{0.36\textwidth}p{0.42\textwidth}@{}}
\toprule
\textbf{Akcja} & \textbf{Wejście} & \textbf{Wywołanie backendu} \\
\midrule
\endhead
\code{create} & \code{name}, \code{game} (lub \code{__custom__} + \code{custom_game}), \code{lobby_location}, \code{players_max}, opcjonalnie \code{description}, \code{communicator_link}, \code{requirements} & \code{POST /rooms} (JSON) \\
\code{join} & \code{room_name} & \code{POST /rooms/{name}/join} \\
\code{leave} & \code{room_name} & \code{POST /rooms/{name}/leave} \\
\code{deleteRoom} & \code{room_name} & \code{DELETE /rooms/{name}} \\
\code{addGame} & \code{name} & \code{POST /games} (JSON \code{{ name }}) \\
\code{deleteGame} & \code{name}, opcjonalnie \code{force="true"} & \code{DELETE /games/{name}?force=true} \\
\bottomrule
\end{longtable}

Szczegóły: w \code{create} wartość \code{__custom__} powoduje użycie \code{custom_game} jako
nazwy gry; \code{players_max} jest parsowane przez \code{parseInt(..., 10)}; pola opcjonalne są
przycinane, a puste łańcuchy zamieniane na \code{null}. W \code{deleteGame} status 409
(aktywne pokoje nadal grają w tę grę) zwraca \code{fail(409, { error, conflict: true, game })},
co pozwala interfejsowi zaproponować wymuszone usunięcie. Wszystkie wywołania backendu idą
przez \code{backendFetch} (rozdział o bibliotece klienta).

\paragraph{Interfejs} (strona ok.\ 1,5 tys.\ linii).
\begin{itemize}[nosep]
  \item \textbf{Lista na żywo.} Singleton \code{roomsStore} łączy się z \code{WS $PUBLIC_WS_URL/ws/rooms}
        (protokół w rozdziale o backendzie). Pokoje wygasłe są odfiltrowywane przy każdym „tyknięciu”.
  \item \textbf{Ping.} Co 15 s strona odpytuje \code{$PUBLIC_BACKEND_URL/healthz} i mierzy czas
        (błąd lub timeout wyświetlany jako \code{ERR}). Szacowane opóźnienie pokoju
        łączy ping API z odległością Haversine między lokalizacją widza a lokalizacją lobby
        (dzieloną przez 100). Trzy koszyki: LOW (do 80 ms, 3 pipsy), MID (do 160 ms, 2), HIGH
        (powyżej 160 ms, 1).
  \item \textbf{Licznik czasu pokoju.} Odliczanie \code{MM:SS} do \code{expires_at}; poniżej 60 s
        ostrzeżenie wizualne.
  \item \textbf{Filtry.} Tekst (nazwa pokoju, gra, etykieta lokalizacji), sloty (2, 4, 8, 16+),
        status ANY/OPEN/FULL, ping ANY/LOW/MID/HIGH.
  \item \textbf{Wybór i panel boczny.} Kliknięcie zaznacza pokój i pokazuje szczegóły; \code{Enter}
        otwiera \code{/rooms/{name}} (wymaga członkostwa lub roli administratora); kliknięcie poza listą
        odznacza.
  \item \textbf{Dialog tworzenia.} Klawisz \code{N} (wymaga logowania). Lista gier lub własna nazwa;
        lokalizacja lobby z listy oraz przycisk „Detect” (\code{navigator.geolocation},
        \code{nearestLobbyLocation()} z \code{$lib/lobbyLocations}); zapisana lokalizacja w
        \code{localStorage} pod kluczem \code{playlink.viewerLocation}.
  \item \textbf{Administrator.} Komponent \code{GameManager} (dodawanie/usuwanie gier z obsługą konfliktu 409)
        oraz bezpośrednie usuwanie dowolnego pokoju z listy.
  \item \textbf{Toasty.} Wyniki akcji formularzy znikają po 4,5 s.
\end{itemize}

\begin{longtable}{@{}p{0.14\textwidth}p{0.75\textwidth}@{}}
\toprule
\textbf{Klawisz} & \textbf{Działanie} \\
\midrule
\endhead
\code{/} & Fokus na polu wyszukiwania \\
\code{R} & Przeładowanie strony \\
\code{N} & Dialog tworzenia (zalogowany) \\
\code{Enter} & Otwórz zaznaczony pokój (członek lub administrator) \\
\code{Esc} & Zamknij dialog / odznacz / \code{history.back()} \\
\bottomrule
\end{longtable}

\fig{d2-rooms-v1.png}{0.9}{Lista pokoi (\code{/rooms}) w stylu Diablo II: nagłówek kolumn Game Name / Game / Players / Ping / Timer, panel filtrów (sloty, status, ping), stan „Awaiting signals…” przy braku aktywnych pokoi oraz pasek podpowiedzi klawiszy. Zrzut z wczesnej wersji, przed dodaniem pozycji Profile w zakładkach.}

\subsection{\texttt{/rooms/[name]} --- szczegóły pokoju}

\paragraph{Load} zwraca: \code{roomName}, \code{roomGame}, \code{roomLocation},
\code{roomLocationLabel}, \code{description}, \code{communicatorLink}, \code{requirements},
\code{token} (surowy JWT, wyłącznie do uwierzytelnienia WebSocket), \code{address},
\code{username}, \code{memberAddresses}, \code{members} (\code{RoomMember[]}), \code{event}
(\code{RoomEventState | null}), \code{isAdmin}, \code{isMember}, \code{isPreview},
\code{createdBy}, \code{isCreator}.
\begin{itemize}[nosep]
  \item Brak sesji lub nieprawidłowy token: \code{redirect(303, '/auth')}.
  \item \code{GET /rooms/{name}} (publiczne); porażka lub \code{null}: przekierowanie na \code{/rooms}.
  \item Użytkownik niebędący członkiem ani administratorem: przekierowanie na \code{/rooms}.
  \item \code{isPreview} oznacza administratora oglądającego pokój, do którego nie należy
        (bez WebSocketu, tylko odczyt).
\end{itemize}

\paragraph{Form actions} (wszystkie wymagają sesji, inaczej \code{fail(401)}):

\begin{longtable}{@{}p{0.14\textwidth}p{0.22\textwidth}p{0.30\textwidth}p{0.26\textwidth}@{}}
\toprule
\textbf{Akcja} & \textbf{Wejście} & \textbf{Wywołanie backendu} & \textbf{Wynik} \\
\midrule
\endhead
\code{closeRoom} & --- & \code{DELETE /rooms/{name}} & \code{{ success, closed: true }} \\
\code{kickMember} & \code{member_address} & \code{POST /rooms/{name}/members/{address}/kick} & \code{{ success, kick }} \\
\code{scheduleEvent} & \code{starts_at}, \code{ends_at} (czas lokalny) & \code{PUT /rooms/{name}/event} (ISO UTC) & \code{{ success, message }} \\
\code{cancelEvent} & --- & \code{DELETE /rooms/{name}/event} & \code{{ success, message }} \\
\code{setRsvp} & \code{status}: \code{present}, \code{absent} lub \code{maybe} & \code{PUT /rooms/{name}/event/rsvp} & \code{{ success, message }} \\
\bottomrule
\end{longtable}

W \code{scheduleEvent} wartości \code{datetime-local} są konwertowane do ISO UTC przez
\code{localInputToIsoUtc()}; nieprawidłowe daty dają \code{fail(400)}. W \code{setRsvp}
status jest walidowany względem zbioru \{\code{present}, \code{absent}, \code{maybe}\}
(inaczej \code{fail(400)}). Błędy backendu są zwracane jako \code{fail(status, { error })}.

\paragraph{Interfejs} (strona ok.\ 1,2 tys.\ linii).
\begin{itemize}[nosep]
  \item \textbf{Czat WebSocket.} Członkowie łączą się przez \code{createChatStore(roomName, token, options)}
        z \code{$lib/chatStore}. Store udostępnia \code{messages}, \code{event}, \code{members},
        \code{owner}, \code{closed}, \code{kicked}. \code{Enter} wysyła, \code{Shift+Enter} wstawia nową
        linię; widok przewija się do dołu przy nowych wiadomościach. Wiadomości są podpisywane kluczem
        z \code{loadSigner()} (\code{sessionStorage}).
  \item \textbf{Lista członków.} Tablica \code{partyMembers} z polami \code{address}, \code{username},
        \code{isMe}, \code{isAdmin}, \code{isOwner}; każdy członek ma mały \code{Sigil}, nazwę (lub skrócony adres
        \code{0x1234...5678}) i odznaki admin/owner.
  \item \textbf{Tryb podglądu administratora.} Bez WebSocketu i pola czatu; stan synchronizowany
        przy unieważnieniu danych (\code{$effect} obserwujący \code{data}).
  \item \textbf{Komponent \code{RoomEvent}.} Planowanie, edycja i anulowanie wydarzenia
        (\code{D2DatePicker} + \code{D2TimePicker}), przyciski RSVP (Present / Absent / Maybe)
        i odliczanie do startu; lista RSVP dopasowuje się do zmian składu.
  \item \textbf{Zamknięcie i wyrzucenie.} \code{closeRoom()} wysyła \code{?/closeRoom} i od razu
        przechodzi do \code{/rooms} (\code{goto()}). \code{kickMember()} wysyła \code{?/kickMember};
        po sukcesie usuwa członka z lokalnej listy i czyści jego RSVP, a przy przekazaniu własności
        aktualizuje \code{ownerAddress}. Gdy WebSocket ogłosi \code{room_closed} lub \code{member_kicked},
        pojawia się modalny \code{SystemDialog} (ton \code{blood}) z 5-sekundowym odliczaniem
        (\code{CLOSE_REDIRECT_SECONDS}) i automatycznym powrotem do \code{/rooms}.
  \item \textbf{Dialogi.} \code{SystemDialog} dla: powiadomienia o wyrzuceniu, o zamknięciu pokoju,
        potwierdzenia wyrzucenia członka i potwierdzenia zamknięcia pokoju.
  \item \textbf{HintBar.} Członkowie: \code{Enter: Transmit}, \code{Shift+Enter: Newline},
        \code{Esc: Leave}; podgląd administratora: \code{Esc: Back to Rooms}.
\end{itemize}

\subsection{\texttt{/test} --- obserwatorium deweloperskie}

Sonda zdrowia backendu na potrzeby rozwoju (bez \code{+page.server.ts}). Po zamontowaniu pobiera
\code{GET $PUBLIC_BACKEND_URL/} i wyświetla surowy JSON, adres backendu, status HTTP i opóźnienie
(ms) wraz z \code{PipMeter} (3 pipsy poniżej 80 ms, 2 poniżej 250 ms, 1 w pozostałych
przypadkach, 0 przy błędzie). Klawisz \code{R} ponawia ping. \code{HintBar}: \code{R: Re-ping},
\code{Esc: Back}; tytuł strony „PlayLink --- Observatory”.

\section{Wczesny prototyp interfejsu (smoke test)}\label{sec:fe-smoke}

Dwa poniższe zrzuty pochodzą z testu dymnego (ang.\ \emph{smoke test}) wcześniejszego prototypu
interfejsu: formularz tworzenia lobby wypełnionymi danymi oraz widok szczegółów pokoju z metadanymi
(opis, link do komunikatora, wymagania). Układ i teksty (np. „Configure Signal”, „Broadcast”)
odbiegają od obecnego UI w stylu Diablo II; pokazują natomiast te same pola, które dziś obsługuje
akcja \code{create} oraz strona \code{/rooms/[name]}: nazwa lobby, gra, liczba slotów,
opis, \code{communicator_link} i \code{requirements}.

\fig{smoke-form-filled.png}{0.62}{Prototyp: formularz tworzenia lobby z wypełnionymi danymi testowymi (nazwa, gra, sloty, opis, link do komunikatora, wymagania).}

\fig{smoke-detail-with-metadata.png}{0.62}{Prototyp: widok szczegółów pokoju z metadanymi (opis, komunikator, wymagania) i pustym czatem.}

%====================================================================
\chapter{Biblioteka klienta i stan}\label{ch:fe-biblioteka}

Rozdział opisuje moduły nie-UI z \code{frontend/src/lib/}: tożsamość i uwierzytelnianie, podpisy
wiadomości czatu, przechowywanie klucza sesji, store'y WebSocket (lista pokoi i czat pokoju),
kontekst podpowiedzi klawiszy, lokalizacje lobby, obsługę regionów, rozwiązywanie adresów URL,
raportowanie błędów klienta oraz hooki serwera. Ramka protokołu WebSocket po stronie serwera jest
opisana w rozdziale o realtime backendu.

\section{Zmienne publiczne i adresy URL}\label{sec:fe-urls}

\begin{longtable}{@{}p{0.27\textwidth}p{0.33\textwidth}p{0.32\textwidth}@{}}
\toprule
\textbf{Zmienna} & \textbf{Użycie} & \textbf{Zastosowanie} \\
\midrule
\endhead
\code{PUBLIC_BACKEND_URL} & \code{urls.ts} (\code{getBackendUrl()}) & Bazowy URL REST w przeglądarce \\
\code{PUBLIC_WS_URL} & \code{urls.ts} (\code{getWsUrl()}) & Bazowy URL WebSocket \\
\bottomrule
\end{longtable}

Moduł \code{lib/urls.ts} normalizuje skonfigurowany adres (przycina białe znaki i końcowy ukośnik)
i dostarcza dwie funkcje:
\begin{itemize}[nosep]
  \item \code{getBackendUrl()}: zwraca \code{PUBLIC_BACKEND_URL}; jeśli brak, a strona działa na
        \code{localhost} lub \code{127.0.0.1}, używa \code{http://<host>:8000}; w innym wypadku loguje
        błąd \code{PUBLIC_BACKEND_URL is not configured} i zwraca pusty łańcuch.
  \item \code{getWsUrl()}: analogicznie dla \code{PUBLIC_WS_URL} (awaryjnie \code{ws://<host>:8000}).
\end{itemize}
Awaryjny adres lokalny istnieje wyłącznie dla zwykłego \code{bun dev} i nie może po cichu kierować
ruchu produkcyjnego na komputer użytkownika, gdy zmienna nie jest ustawiona.
\code{auth.ts} używa \code{getBackendUrl()}, a \code{roomsStore.ts} i \code{chatStore.ts}
używają \code{getWsUrl()}.

\section{Tożsamość i uwierzytelnianie (\texttt{auth.ts})}\label{sec:fe-auth}

Trójfazowy przepływ oparty na mnemoniku BIP39 i podpisach EIP-191 (\code{personal_sign}), ethers v6.

\begin{longtable}{@{}p{0.36\textwidth}p{0.58\textwidth}@{}}
\toprule
\textbf{Funkcja} & \textbf{Opis} \\
\midrule
\endhead
\code{getIdentityFromMnemonic(phrase): HDNodeWallet} & Derywuje \code{HDNodeWallet} z 12-wyrazowej frazy BIP39; \code{wallet.address} (adres z sumą kontrolną) jest tożsamością użytkownika. Błąd: \code{Error('Invalid mnemonic phrase')}. \\
\code{generateMnemonic(): string} & Nowa fraza 12 słów z 16 bajtów entropii (\code{Mnemonic.fromEntropy(ethers.randomBytes(16))}). \\
\code{signMessage(identity, message): Promise<string>} & Podpis EIP-191; zwraca \code{0x}-hex (65 bajtów: \code{r || s || v}). \\
\code{authenticate(phrase): Promise<{ token, address, username }>} & Pełny handshake z backendem (poniżej). \\
\bottomrule
\end{longtable}

\paragraph{Faza 1 --- nonce.} \code{POST {PUBLIC_BACKEND_URL}/auth/request-nonce?address=<adres>}
bez uwierzytelnienia (limit żądań wg IP/adresu). Odpowiedź 200: \code{{ nonce }} (UUID v4).
Błędy: 429 czyta pole \code{error.error}, pozostałe \code{error.detail}.

\paragraph{Faza 2 --- podpis lokalny.} Podpisywana wiadomość (pojedynczy znak \code{\n} U+000A,
bez znaku końcowego): \code{Sign in to Playlink\nNonce: <nonce>}. Klucz prywatny jest
zapisywany w \code{sessionStorage} przez \code{saveSigningKey(identity.privateKey)} zaraz po derywacji.

\paragraph{Faza 3 --- weryfikacja.} \code{POST {PUBLIC_BACKEND_URL}/auth/verify},
\code{Content-Type: application/json}, ciało \code{{ address, nonce, signature }}.
Odpowiedź 200: \code{{ token, username }}; frontend przekazuje token do akcji \code{?/login}
(sekcja o modelu BFF). Błędy: \code{error.detail}.

\section{Podpisy wiadomości czatu (\texttt{signing.ts})}\label{sec:fe-signing}

Każda wiadomość czatu jest podpisywana (EIP-191), aby backend mógł kryptograficznie udowodnić
autorstwo (zgłoszenie \#59). Kanoniczny ciąg musi być zgodny co do bajta z funkcją backendu
\code{_chat_signing_message}.

\begin{itemize}[nosep]
  \item \code{ChatSigner}: interfejs z metodą \code{signMessage(message): Promise<string>}; spełniają go
        \code{ethers.HDNodeWallet} i \code{ethers.Wallet}.
  \item \code{buildChatSigningMessage(room, content, sentAt)} składa tekst:
\end{itemize}

\begin{lstlisting}
PlayLink signed chat message
room=<room>
sent_at=<sentAt>
content=<content>
\end{lstlisting}

Podział linii to pojedynczy znak U+000A; \code{room} to nazwa pokoju widoczna dla użytkownika
(nie zakodowana URL), \code{sentAt} to \code{new Date().toISOString()}, a \code{content} to
przycięta treść. \code{signChatMessage(signer, room, content, sentAt)} buduje ciąg i podpisuje go,
zwracając \code{0x}-hex.

\section{Przechowywanie klucza (\texttt{signingKey.ts})}\label{sec:fe-signingkey}

Klucz prywatny ECDSA wyprowadzony z BIP39 jest trzymany w \code{sessionStorage} (klucz
\code{playlink.sk}) na czas życia karty; nigdy w ciasteczku ani \code{localStorage}. Model
bezpieczeństwa: klucz jest narażony na XSS --- jak każdy sekret na stronie w aplikacji
samodzielnie przechowującej klucze, bez rozszerzenia przeglądarki.

\begin{itemize}[nosep]
  \item \code{saveSigningKey(privateKey)}: zapis (no-op poza przeglądarką, błąd logowany).
  \item \code{loadSigner(): Wallet | null}: zwraca \code{ethers.Wallet} lub \code{null}, gdy brak klucza
        albo jest nieprawidłowy (nieprawidłowy klucz jest po cichu usuwany).
  \item \code{clearSigningKey()}: usuwa klucz (przy wylogowaniu).
\end{itemize}

\section{Store listy pokoi (\texttt{roomsStore.ts})}\label{sec:fe-roomsstore}

Store oparty o WebSocket, utrzymujący listę aktywnych pokoi.

\paragraph{Typ \texttt{RoomSummary}:} \code{name}, \code{game}, \code{lobby_location},
\code{players_active}, \code{players_max}, \code{member_addresses: string[]},
\code{description}, \code{communicator_link}, \code{requirements} (trzy ostatnie
\code{string | null}), \code{expires_at}. Funkcje \code{isNullableString} i \code{isRoomSummary}
walidują nieufne dane; każda ramka z WebSocketu przechodzi przez tę walidację.

\paragraph{\texttt{roomsStore}} (singleton \code{Readable<RoomSummary[]> & { destroy() }}):
\begin{itemize}[nosep]
  \item \textbf{Adres:} \code{{PUBLIC_WS_URL}/ws/rooms}, bez uwierzytelniania (lista publiczna).
  \item \textbf{Format ramki:} goła tablica JSON obiektów \code{RoomSummary}, bez koperty
        \code{{ type }}; walidacja \code{Array.isArray(data) && data.every(isRoomSummary)}, potem \code{set(data)}.
  \item \textbf{Ponowne łączenie:} wykładniczy backoff z jitterem: baza 1000 ms, maksimum 30\,000 ms,
        \code{delay = min(30000, 1000 * 2^attempts) + random * (0.5 * wykładniczy)}; licznik zerowany po
        udanym otwarciu.
  \item \textbf{\code{destroy()}:} zamyka gniazdo, kasuje zaplanowane ponowienia i ustawia
        \code{isTornDown = true}.
\end{itemize}

\section{Store czatu pokoju (\texttt{chatStore.ts})}\label{sec:fe-chatstore}

Fabryka \code{createChatStore(roomName, token, options?)} zwraca osobny \code{ChatStore}
dla każdego pokoju.

\subsection{Typy}

\begin{longtable}{@{}p{0.22\textwidth}p{0.72\textwidth}@{}}
\toprule
\textbf{Typ} & \textbf{Pola} \\
\midrule
\endhead
\code{ChatMessage} & \code{id}, \code{sender_address}, \code{sender_username}, \code{content}, \code{created_at}; opcjonalnie \code{kind} (\code{'user' | 'system'}), \code{signature} (EIP-191 lub \code{null} dla wiadomości historycznych), \code{sent_at} (dokładny podpisany znacznik klienta), \code{verified} (serwer odzyskał prawidłowego sygnatariusza) \\
\code{RoomMember} & \code{address}, \code{username}, \code{is_admin} \\
\code{RsvpStatus} & \code{'present' | 'absent' | 'maybe'} \\
\code{RsvpEntry} & \code{address}, \code{username}, \code{status}, \code{updated_at} \\
\code{RoomEventState} & \code{starts_at}, \code{ends_at}, \code{created_by}, \code{created_at}, \code{updated_at}, \code{rsvps: RsvpEntry[]} \\
\bottomrule
\end{longtable}

Walidatory \code{isChatMessage}, \code{isRoomMember}, \code{isRsvpStatus}, \code{isRsvpEntry},
\code{isRoomEventState} sprawdzają strukturę w czasie działania.

\subsection{Ramki obsługiwane przez \texttt{parseFrame}}

\code{parseFrame(raw): ChatFrame | null} parsuje i waliduje surowy komunikat; przy błędnym JSON-ie
lub nieznanym typie zwraca \code{null}.

\begin{longtable}{@{}p{0.17\textwidth}p{0.30\textwidth}p{0.42\textwidth}@{}}
\toprule
\textbf{\code{type}} & \textbf{Dodatkowe pola} & \textbf{Działanie store'a} \\
\midrule
\endhead
\code{history} & \code{messages} & \code{messages.set(...)} --- zastępuje całość \\
\code{message} & \code{message} & dopisuje wiadomość do \code{messages} \\
\code{event_update} & \code{event} (lub \code{null}) & \code{event.set(...)} \\
\code{rsvp_update} & \code{rsvp} & scala RSVP w \code{event} (dopasowanie adresu bez rozróżniania wielkości liter) \\
\code{roster_update} & \code{members} & \code{members.set(...)} --- zastępuje całość \\
\code{room_closed} & \code{room} & \code{closed.set(true)}, \code{isTornDown = true}, zamknięcie gniazda \\
\code{member_kicked} & \code{member_address}, \code{member_username}, \code{created_by}, \code{ownership_transferred} & \code{owner.set(created_by)}; gdy \code{member_address} to \code{options.currentAddress}: \code{kicked.set(true)}, \code{isTornDown = true} \\
\code{error} & \code{detail} & \code{console.warn('Chat message rejected:', detail)} \\
\bottomrule
\end{longtable}

\subsection{Interfejs \texttt{ChatStore} i opcje}

\begin{longtable}{@{}p{0.17\textwidth}p{0.20\textwidth}p{0.55\textwidth}@{}}
\toprule
\textbf{Pole} & \textbf{Wartość początkowa} & \textbf{Opis} \\
\midrule
\endhead
\code{messages} & \code{[]} & Historia czatu pokoju na żywo \\
\code{event} & \code{initialEvent ?? null} & Bieżące wydarzenie (lub \code{null}) \\
\code{members} & \code{initialMembers ?? []} & Skład pokoju, aktualizowany przy dołączeniu/wyjściu \\
\code{closed} & \code{false} & \code{true} po zamknięciu pokoju (\code{room_closed}) \\
\code{owner} & \code{initialOwner ?? ''} & Adres właściciela; zmienia się po wyrzuceniu przez twórcę \\
\code{kicked} & \code{false} & \code{true}, gdy ten klient został usunięty z pokoju \\
\code{send(content)} & --- & Wysyła wiadomość (poniżej) \\
\code{destroy()} & --- & Zamyka gniazdo i anuluje ponowienia \\
\bottomrule
\end{longtable}

Opcje \code{CreateChatStoreOptions}: \code{initialEvent}, \code{initialMembers}, \code{initialOwner}
(dane z SSR eliminujące migotanie pustej treści do czasu pierwszej ramki WebSocket),
\code{currentAddress} (wykrywanie własnego wyrzucenia) oraz \code{signer}
(\code{ChatSigner | null}; bez niego wiadomości idą niepodpisane).

\subsection{Połączenie i wysyłanie}

\begin{itemize}[nosep]
  \item \textbf{URL:} \code{{PUBLIC_WS_URL}/ws/rooms/{encodeURIComponent(roomName)}/chat?token={encodeURIComponent(token)}}.
        JWT idzie w parametrze zapytania, ponieważ przeglądarki nie pozwalają ustawiać nagłówków WebSocket.
  \item \textbf{Ponowne łączenie:} jak w \code{roomsStore} (1000 ms, maks.\ 30\,000 ms, jitter 50\%);
        wstrzymane po \code{closed}, \code{kicked} oraz kodach zamknięcia 4403 i 4409.
  \item \textbf{\code{send(content)}:} (1) przycina treść; nic nie robi, gdy pusta lub gniazdo nie jest \code{OPEN};
        (2) przy ustawionym \code{signer} generuje \code{sentAt}, wywołuje \code{signChatMessage} i wysyła
        \code{{"content", "sent_at", "signature"}}; (3) przy błędzie podpisu wraca do wiadomości
        niepodpisanej \code{{"content"}} i loguje \code{console.error}; (4) bez signera wysyła
        \code{{"content"}}.
\end{itemize}

\begin{longtable}{@{}p{0.16\textwidth}p{0.75\textwidth}@{}}
\toprule
\textbf{Kod zamknięcia} & \textbf{Zachowanie} \\
\midrule
\endhead
4403 & Nie jesteś członkiem pokoju: \code{kicked.set(true)}, bez ponowień \\
4409 & Nieznany pokój: \code{kicked.set(true)}, bez ponowień \\
pozostałe & Planowane jest ponowne połączenie (o ile store nie został zdemontowany) \\
\bottomrule
\end{longtable}

\begin{lstlisting}
const store = createChatStore(roomName, token, { signer, currentAddress });
// ... store.messages, store.send(...) ...
store.destroy();   // zamyka gniazdo, anuluje ponowienia
\end{lstlisting}

\section{Kontekst podpowiedzi klawiszy (\texttt{hintsContext.svelte.ts})}\label{sec:fe-hints}

Reaktywny kontekst (runes \code{$state}) zasilający pasek podpowiedzi. Typ
\code{HintTone}: \code{gold}, \code{green}, \code{amber}, \code{red}, \code{stone}, \code{blue}.
\code{HintEntry}: \code{key}, \code{label}, opcjonalnie \code{tone} i \code{onclick}.
Klasa \code{HintsState} ma reaktywne pole \code{hints} oraz metody \code{set(entries)} i \code{clear()}.
\code{provideHints(initial?)} tworzy stan i rejestruje go w kontekście komponentu pod kluczem
\code{Symbol('hints')} (wywoływane raz w layoucie lub stronie); \code{getHintsState()} odczytuje stan
z kontekstu (\code{undefined} poza poddrzewem, które wywołało \code{provideHints}). Layout główny
tworzy domyślny stan (\code{Esc: Back}, \code{?: Help}), a strony nadpisują go własnymi skrótami.

\section{Lokalizacje lobby (\texttt{lobbyLocations.ts})}\label{sec:fe-lobby}

Typ \code{LobbyLocation} to \code{{ code, label, lat, lon }}. Stała
\code{FALLBACK_LOBBY_LOCATION = 'eu-central'}, a \code{FALLBACK_LOBBY_LOCATIONS} zawiera 10
lokalizacji awaryjnych, używanych, gdy backend nie zwróci listy:

\begin{longtable}{@{}p{0.2\textwidth}p{0.34\textwidth}p{0.12\textwidth}p{0.12\textwidth}@{}}
\toprule
\textbf{Kod} & \textbf{Etykieta} & \textbf{Szer.} & \textbf{Dł.} \\
\midrule
\endhead
\code{na-east} & North America East & 39,0 & $-77,0$ \\
\code{na-west} & North America West & 37,8 & $-122,4$ \\
\code{eu-west} & Europe West & 51,5 & $-0,1$ \\
\code{eu-central} & Europe Central & 50,1 & 8,7 \\
\code{eu-north} & Europe North & 59,3 & 18,1 \\
\code{sa-east} & South America East & $-23,6$ & $-46,6$ \\
\code{asia-east} & Asia East & 35,7 & 139,7 \\
\code{asia-south} & Asia South & 1,3 & 103,8 \\
\code{oceania} & Oceania & $-33,9$ & 151,2 \\
\code{africa-south} & Africa South & $-26,2$ & 28,0 \\
\bottomrule
\end{longtable}

Funkcje pomocnicze: \code{isLobbyLocation(value)} (strażnik typu), \code{lobbyLocationLabel(locations, code)}
(etykieta lub sam kod, gdy brak dopasowania), \code{distanceKm(a, b)} (odległość wielkokołowa wzorem
Haversine'a, promień Ziemi 6371 km, wynik w km) oraz \code{nearestLobbyLocation(locations, coords)}
(lokalizacja o najmniejszej odległości lub \code{null} dla pustej tablicy).

\section{Hooki serwera i typy globalne}\label{sec:fe-hooks}

\paragraph{\texttt{handle}} (\code{hooks.server.ts}) uruchamia się przy każdym żądaniu serwerowym:
czyta ciasteczko \code{session}, dekoduje JWT (\code{sub}, \code{exp}, \code{is_admin}), sprawdza
wygaśnięcie (\code{exp} w sekundach, porównywane w ms) i ustawia
\code{event.locals.user = { address: sub, isAdmin: is_admin === true }}. Bez ciasteczka
\code{locals.user} pozostaje \code{undefined}; wygasły lub uszkodzony token powoduje usunięcie
ciasteczka (\code{path: '/'}).

\paragraph{\texttt{handleError}} (serwer) zapisuje każdy nieobsłużony błąd jako wiersz logu
\code{server_error} (poziom ERROR; pola \code{path}, \code{status}, \code{error}, \code{stack})
i zwraca klientowi neutralny komunikat \code{{ message: 'Internal Error' }}.

\paragraph{\texttt{app.d.ts}} rozszerza \code{App.Locals} o opcjonalne
\code{user?: { address: string; isAdmin: boolean }}, co zapewnia typowanie \code{event.locals.user}
w hookach, loaderach i akcjach.

\section{Fetch po stronie serwera (\texttt{lib/server/backend.ts})}\label{sec:fe-backend}

Loadery i form actions nie wołają \code{fetch} bezpośrednio, tylko przez dwie funkcje:

\begin{itemize}[nosep]
  \item \code{backendBase()}: wewnętrzny adres backendu: \code{BACKEND_INTERNAL_URL} (w compose nazwa
        usługi), w przeciwnym razie \code{PUBLIC_BACKEND_URL}, ostatecznie \code{http://localhost:8000}
        (żeby \code{bun dev} działał poza Dockerem).
  \item \code{backendFetch(event, path, init)}: wywołuje \code{fetch(backendBase() + path)} i
        \textbf{przekazuje prawdziwy adres IP gracza} w nagłówku \code{CF-Connecting-IP} (jeśli
        żądanie go zawiera). Backend ufa temu nagłówkowi tylko od kontenera frontendu, więc bez
        przekazania wszystkie wywołania SSR dzieliłyby jeden kubełek limitu żądań. Nagłówki są
        trzymane jako zwykły obiekt, aby istniejące testy i miejsca wywołań zachowały dokładną
        pisownię \code{Authorization}.
\end{itemize}

\section{Regiony w interfejsie}\label{sec:fe-regiony}

Platforma działa w dwóch regionach, EU i US (architektura i wdrożenie --- patrz część o
architekturze i operacjach). Interfejs odpowiada za: wybór regionu przy pierwszej wizycie,
zapamiętanie wyboru, propozycję zmiany na szybszy region (Stay/Switch) i automatyczny failover,
gdy zapisany region nie odpowiada. Zasada nadrzędna: \emph{szybszy region jest tylko
propozycją; automatyczne przełączenie następuje wyłącznie wtedy, gdy zapisany region milczy}.

\subsection{Konfiguracja regionów (\texttt{lib/regions.ts})}

\code{REGIONS} zawiera dwa wpisy \code{RegionConfig} \code{{ id, label, frontend, api, ws, healthPath }}:
\code{eu} (etykieta „Europa”) i \code{us} (etykieta „USA”), z \code{healthPath = '/healthz'}.
Adresy pochodzą ze zmiennych \code{PUBLIC_REGION_{EU,US}_{FRONTEND,API,WS}} (przycięte, bez końcowego
ukośnika). W trybie deweloperskim brakujące wartości mają awaryjne adresy lokalne
(EU: 3000/8000, US: 3001/8001); w produkcji brak zmiennej daje pusty łańcuch, aby nie kierować
ruchu na \code{localhost}. Dodatkowo:

\begin{itemize}[nosep]
  \item \code{ENTRY_HOST} (z \code{PUBLIC_ENTRY_HOST}) i \code{isEntryHost(host)}: czy dany host
        serwuje wybór regionu;
  \item \code{CURRENT_REGION}: region tej instancji (\code{PUBLIC_REGION === 'us'} daje \code{us},
        w przeciwnym razie \code{eu});
  \item \code{otherRegion(id)}: drugi region.
\end{itemize}

\subsection{Pomiar opóźnienia i decyzja o propozycji}

\begin{itemize}[nosep]
  \item \code{probeRegion(region, samples = 2)}: dwukrotnie wykonuje
        \code{fetch({api}/healthz, { cache: 'no-store' })} z limitem czasu 2500 ms
        (\code{AbortController}); zwraca najlepszy udany czas (\code{ProbeResult}: \code{ok}, \code{ms}).
        Timeout, błąd sieci lub odpowiedź inna niż OK oznaczają „nie kandydat”. Region bez
        skonfigurowanego \code{api} jest od razu nieosiągalny.
  \item \code{probeAll(regions)}: sonduje wszystkie regiony równolegle.
  \item \code{liveRegionsSorted(results)}: żywe regiony posortowane od najszybszego.
  \item \code{isMeaningfullyFaster(candidateMs, baselineMs, { minDeltaMs, minRatio })}: czy kandydat
        jest \emph{wyraźnie} szybszy od bazy. Domyślne progi: różnica co najmniej 25 ms \emph{lub}
        czas nie większy niż 0,85 bazy. Dla \code{localhost} i \code{127.0.0.1} progi są łagodniejsze
        (5 ms lub 0,98), aby propozycję dało się pokazać w demonstracjach laboratoryjnych. Brak
        pomiaru po którejkolwiek stronie daje \code{false}.
\end{itemize}

\subsection{Ciasteczko wyboru i przekierowanie}

\begin{itemize}[nosep]
  \item \textbf{Ciasteczko \code{playlink.region}} (wartość \code{eu} lub \code{us}), ustawiane po stronie klienta
        przez \code{saveRegion(id)}: \code{Path=/}, \code{Max-Age=31536000} (1 rok), \code{SameSite=Lax},
        \code{Secure} na \code{https:}. \code{readSavedRegion()} odczytuje je i akceptuje wyłącznie
        \code{eu} i \code{us}.
  \item \textbf{\code{redirectToRegion(id, regions)}:} przekierowuje przez \code{window.location.replace}
        na frontend regionu, \emph{zachowując ścieżkę i zapytanie}. Zwraca \code{false}, gdy cel jest tym
        samym hostem (ochrona przed pętlą przekierowań). Origin jest „przypięty”: konstruowany z
        konfiguracji regionu, a z bieżącego adresu przenoszone są tylko \code{pathname} (z początkowymi
        ukośnikami zredukowanymi do jednego) i \code{search}, więc ścieżka typu \code{//evil.example}
        nie może zmienić hosta.
\end{itemize}

\subsection{Wybór regionu (\texttt{RegionChooser.svelte})}

Komponent renderowany na hoście wejściowym (\code{isEntryHost}). Pokazuje herb, tytuł, panel
„REGION SELECT” ze statusem, listę regionów z opóźnieniem (\code{NN ms} lub \code{offline}) i
przyciskiem „Open” oraz notatki dla użytkownika. Stan: \code{probing}, \code{redirecting},
\code{propose}, \code{failed}. Logika przy zamontowaniu, w kolejności:

\begin{enumerate}[nosep]
  \item \textbf{Wymuszenie przez URL.} Parametr \code{?region=eu|us} od razu wybiera dany region
        (zapis ciasteczka i przekierowanie).
  \item \textbf{Pomiar.} \code{probeAll} mierzy oba regiony; wyniki żywe są sortowane.
  \item \textbf{Zapisany region nie odpowiada} (jest w ciasteczku, ale sonda nieudana): failover do
        najszybszego żywego regionu („Saved region \ldots is offline --- opening \ldots”). Gdy żaden nie
        odpowiada: stan \code{failed}, komunikat „Saved region is offline and no fallback is responding.”
  \item \textbf{Zapisany region żyje.} Pozostaje się na nim, chyba że drugi żywy region jest
        wyraźnie szybszy (\code{isMeaningfullyFaster}); wtedy \emph{tylko proponuje się zmianę} (stan
        \code{propose}) z komunikatem w rodzaju „Europa looks faster (NN ms vs MM ms). Stay on \ldots,
        or switch?” i dwoma przyciskami: „Stay on \{zapisany\}” (\code{keepSaved}) i „Switch to
        \{proponowany\}” (\code{acceptProposal}). Każdy wybór zapisuje ciasteczko i przekierowuje.
  \item \textbf{Pierwsza wizyta} (brak ciasteczka): otwierany jest najszybszy żywy region; gdy żaden nie
        odpowiada --- stan \code{failed} z komunikatem „Neither region is responding” i bezpośrednimi
        linkami do obu frontendów.
  \item \textbf{Pomocnik laboratoryjny} \code{?propose=1}: wymusza ekran Stay/Switch nawet, gdy oba
        regiony są równie szybkie (przy braku ciasteczka zapisany region jest „udawany” jako wolniejszy
        z dwóch), by dało się wykonać zrzuty ekranu do dziennika.
\end{enumerate}

Jeśli wybrany region ma ten sam host co bieżący, \code{select} kończy się stanem \code{failed}
(„You are already on the \ldots host. Open the lobby manually.”). Notatki na ekranie przypominają, że:
pierwsza wizyta otwiera szybszy region, wybór jest zapamiętywany, szybszy drugi region to tylko
propozycja, automatyczne przełączenie następuje wyłącznie przy braku odpowiedzi zapisanego regionu, a
\emph{każdy region ma osobną sesję} --- po przełączeniu trzeba zalogować się tą samą frazą 12 słów.

Schemat całego procesu wyboru regionu jest w części II, w opisie filaru F2.

\subsection{Baner propozycji (\texttt{RegionProposalBanner.svelte})}

Na hoście regionalnym layout osadza baner, który \emph{po fakcie} proponuje zmianę regionu, gdy
użytkownik już pracuje w jednym z nich. Propsy: \code{currentRegion} (\code{RegionId}) i \code{other}
(\code{RegionConfig | undefined}). Przy zamontowaniu (jeśli drugi region ma skonfigurowane \code{api}
i baner nie został wcześniej odrzucony) uruchamia \code{probeAll()}; baner pojawia się tylko, gdy
\emph{oba} regiony odpowiedziały i drugi jest wyraźnie szybszy (\code{isMeaningfullyFaster}).
Treść: „\{inny\} looks faster (NN ms vs MM ms). Saved region stays put unless you switch.” Akcje:
„Keep current” (zapamiętuje odrzucenie w \code{sessionStorage} pod kluczem
\code{playlink.regionProposalDismissed} wartością bieżącego regionu, więc baner nie wraca w tej karcie)
oraz „Switch to \{inny\}” (link do tej samej ścieżki i zapytania w drugim regionie; kliknięcie wywołuje
\code{saveRegion(other.id)}). Obok odznaki aktywnego regionu layout zawiera także stały link
ręcznego przełączenia „$\rightarrow$ \{druga etykieta\}”.

\subsection{Podsumowanie zachowań regionalnych}

\begin{longtable}{@{}p{0.30\textwidth}p{0.62\textwidth}@{}}
\toprule
\textbf{Sytuacja} & \textbf{Zachowanie interfejsu} \\
\midrule
\endhead
Pierwsza wizyta na hoście wejściowym & Pomiar obu regionów, otwarcie najszybszego żywego, zapis ciasteczka \\
Zapisany region żyje, drugi wyraźnie szybszy & Propozycja Stay/Switch; domyślnie pozostaje zapisany \\
Zapisany region milczy (sonda nieudana) & Failover do najszybszego żywego regionu (jedyny przypadek automatycznej zmiany) \\
Żaden region nie odpowiada & Komunikat o błędzie i bezpośrednie linki do frontendów \\
Użytkownik już w regionie, drugi szybszy & Baner z „Keep current” / „Switch to \ldots”; odrzucenie na czas karty \\
Przełączenie regionu & Osobna sesja (ciasteczko host-only): ponowne logowanie tą samą frazą \\
\bottomrule
\end{longtable}

\section{Sondy zdrowia i raportowanie błędów klienta}\label{sec:fe-bledy}

\subsection{\texttt{GET /healthz}}

Endpoint serwerowy frontendu (\code{routes/healthz/+server.ts}) zwraca
\code{{"status":"ok"}} z nagłówkami \code{content-type: application/json} i
\code{cache-control: no-store}. Służy sondom zdrowia infrastruktury. Uwaga: pomiary opóźnienia
regionów (\code{probeRegion}) oraz ping na stronie \code{/rooms} wołają \code{/healthz}
\emph{backendu} (API), nie frontendu.

\subsection{Raportowanie błędów przeglądarki}

\begin{itemize}[nosep]
  \item \textbf{Zbieranie (\code{hooks.client.ts}).} Rejestruje nasłuch zdarzeń \code{error} i
        \code{unhandledrejection} oraz \code{handleError} SvelteKit; wszystkie przekazują błąd do
        \code{reportClientError()} (użytkownikowi pokazywany jest neutralny komunikat
        \code{Something went wrong}).
  \item \textbf{Wysyłka (\code{lib/clientErrors.ts}).} \code{reportClientError(error)} wysyła
        \code{navigator.sendBeacon('/api/client-errors', ...)} z JSON-em \code{{ path, message, stack }};
        działa „po najlepszych staraniach”, nie więcej niż \textbf{5 razy na jedno załadowanie strony}.
  \item \textbf{Odbiór (\code{routes/api/client-errors/+server.ts}).} \code{POST} z limitem
        \textbf{20 żądań na minutę} liczonym względem równorzędnego peera połączenia
        (\code{getClientAddress()}), nigdy nagłówka, który klient mógłby sfałszować; za tunelem peerem
        jest konektor, więc limit jest współdzielony przez wszystkich odwiedzających. Przekroczenie:
        429. Ciało niebędące obiektem JSON: 400. Poprawne zgłoszenie jest zapisywane jako wiersz logu
        \code{client_error} (poziom ERROR) z polami \code{path} (do 200 znaków), \code{error} (300),
        \code{stack} (1500) i \code{user_agent} (200) i zwraca 204.
  \item \textbf{Log (\code{lib/server/log.ts}).} \code{logLine(level, msg, fields)} zapisuje jedną linię
        JSON na wpis (\code{_time}, \code{_msg}, \code{level}, \code{logger: "frontend"} oraz pola) na
        stdout (ERROR na stderr). Agent regionu przesyła stdout do VictoriaLogs --- tak błędy
        przeglądarki i serwera frontendu trafiają do wspólnej obserwowalności (szczegóły w części
        o operacjach).
\end{itemize}

%====================================================================
\chapter{Komponenty interfejsu}\label{ch:fe-komponenty}

Frontend zawiera trzy komponenty domenowe (\code{GameManager}, \code{MnemonicInput},
\code{RoomEvent}), dwa komponenty regionalne (\code{RegionChooser}, \code{RegionProposalBanner};
opisane w rozdziale o bibliotece klienta) oraz zestaw 18 widżetów w stylu Diablo II (UI kit
\code{chrome}). Wszystkie korzystają z konwencji Svelte 5 runes: \code{$props()} dla propsów,
\code{$bindable()} dla stanu dwukierunkowego, \code{$state}/\code{$derived}/\code{$effect}
dla stanu lokalnego oraz \code{{@render children?.()}} dla slotów (snippetów).

Wspólne tokeny projektowe (zmienne CSS: kolory, typografia, odstępy) są w
\code{frontend/src/lib/styles/tokens.css}, a atomowe klasy (fazowanie, wytrawianie, kapitaliki) w
\code{styles/chrome.css}. \code{lib/global.css} zawiera reset CSS, deklaracje fontów (Exocet jako
display, Cinzel jako mono) i kursory w stylu Diablo II. Pliki te importuje layout główny;
poszczególne komponenty nie importują ich bezpośrednio.

\section{Komponenty domenowe}\label{sec:fe-domenowe}

\subsection{GameManager}

Dialog dla administratora do tworzenia i usuwania kategorii gier. Wysyła akcje \code{?/addGame}
i \code{?/deleteGame} przez \code{fetch} + \code{deserialize}.

\begin{longtable}{@{}p{0.13\textwidth}p{0.17\textwidth}p{0.14\textwidth}p{0.50\textwidth}@{}}
\toprule
\textbf{Prop} & \textbf{Typ} & \textbf{Bindable} & \textbf{Opis} \\
\midrule
\endhead
\code{open} & \code{boolean} & tak & Widoczność dialogu \\
\code{games} & \code{string[]} & nie & Bieżąca lista kategorii gier \\
\code{onclose} & \code{() => void} & nie & Wywoływane przy zamknięciu \\
\bottomrule
\end{longtable}

Komponent opakowuje \code{SystemDialog} (\code{tone="gold"}, \code{modal}, \code{width="520px"}).
Stan wewnętrzny: pole nazwy nowej gry, komunikat \code{feedback} (\code{'success' | 'error'}), flaga
\code{busy} i obiekt \code{conflict}. Przepływ konfliktu przy usuwaniu: (1) \code{deleteGame(game)} wysyła
\code{?/deleteGame} z \code{name}; (2) odpowiedź \code{{ type: 'failure', data: { conflict: true, error } }}
pokazuje „Force Delete” i link „Cancel”; (3) „Force Delete” wywołuje \code{deleteGame(game, true)},
dodając \code{force: 'true'}; (4) po sukcesie \code{invalidateAll()} odświeża dane trasy. Klawisz
\code{Enter} w polu nazwy wywołuje \code{addGame()} (\code{onAddKey}). Implementacje akcji
opisano w rozdziale o trasach (\code{rooms/+page.server.ts}).

\subsection{MnemonicInput}

Siatka wprowadzania 12 słów BIP39: trzy kolumny po cztery pola z numeracją i walidacją słowa
względem słownika BIP39 (\code{ethers.wordlists.en.getWordIndex()}). Słowa są trzymane jako
\code{string[]} (12 pozycji) i synchronizowane z bindowalnym \code{value} (łańcuch rozdzielany spacjami).

\begin{longtable}{@{}p{0.13\textwidth}p{0.17\textwidth}p{0.14\textwidth}p{0.50\textwidth}@{}}
\toprule
\textbf{Prop} & \textbf{Typ} & \textbf{Bindable} & \textbf{Opis} \\
\midrule
\endhead
\code{value} & \code{string} & tak & Fraza 12 słów rozdzielonych spacjami \\
\bottomrule
\end{longtable}

\paragraph{Synchronizacja dwukierunkowa.} Rodzic $\rightarrow$ dziecko: \code{$effect} obserwuje
\code{value}; gdy staje się 12-wyrazowym łańcuchem, aktualizuje \code{words} (tylko jeśli wartość
zewnętrzna różni się od stanu wewnętrznego); wyczyszczenie \code{value} zeruje wszystkie pola.
Dziecko $\rightarrow$ rodzic: \code{$effect} łączy \code{words} spacją i zapisuje do \code{value}.

\paragraph{Nawigacja.} \code{Space}, \code{Enter}, \code{Tab} --- następne pole; \code{Backspace} w pustym
polu --- poprzednie; wklejenie dzieli tekst po białych znakach, wypełnia do 12 pól, zamienia na małe litery
i usuwa znaki niealfabetyczne.

\paragraph{Walidacja pola.} Niepuste słowo spoza słownika dostaje klasę \code{invalid} (czerwony odcień);
poprawne --- \code{valid-filled} (zielony odcień) i znacznik \code{▸}. Używany na trasie \code{/auth}.

\subsection{RoomEvent}

Panel planowania i RSVP wydarzenia pokoju: odliczanie, lista RSVP oraz (dla twórcy) sterowanie
edycją i anulowaniem, a (dla członków) przyciski RSVP.

\begin{longtable}{@{}p{0.17\textwidth}p{0.27\textwidth}p{0.10\textwidth}p{0.38\textwidth}@{}}
\toprule
\textbf{Prop} & \textbf{Typ} & \textbf{Domyśl.} & \textbf{Opis} \\
\midrule
\endhead
\code{event} & \code{RoomEventState | null} & --- & Dane wydarzenia lub \code{null}, gdy niezaplanowane \\
\code{isCreator} & \code{boolean} & --- & Widz jest właścicielem pokoju \\
\code{isMember} & \code{boolean} & --- & Widz jest członkiem pokoju \\
\code{viewerAddress} & \code{string} & --- & Adres Ethereum widza (wewnętrznie małymi literami) \\
\code{members} & \code{{ address: string; username?: string }[]} & --- & Wszyscy członkowie pokoju \\
\code{formError} & \code{string | null} & \code{null} & Błąd serwera z formularza rodzica \\
\bottomrule
\end{longtable}

\paragraph{Stany.} (1) Brak wydarzenia, widz jest twórcą: przycisk „Schedule Event” przełącza tryb edycji.
(2) Brak wydarzenia, widz nie jest twórcą: komunikat „No event scheduled yet”. (3) Wydarzenie zaplanowane,
bez edycji: czas startu i końca (\code{Intl.DateTimeFormat}), odliczanie, sterowanie RSVP (dla członków)
i lista. (4) Edycja: \code{D2DatePicker} + \code{D2TimePicker} dla startu i końca; po zatwierdzeniu walidacja
lokalna i wysyłka do \code{?/scheduleEvent}.

\paragraph{Odliczanie} (odświeżane co 1 s przez \code{setInterval}): przed startem
\code{STARTS IN [d HH:MM:SS]}; w trakcie \code{IN PROGRESS --- [HH:MM:SS] LEFT}; po końcu \code{ENDED}.

\paragraph{RSVP.} Trzy przyciski: Present, Absent, Maybe; każdy wysyła \code{?/setRsvp} z \code{status}.
Lista scala \code{event.rsvps} z \code{members} (adresy bez wpisu w \code{rsvps} mają „No Response”).

\paragraph{Walidacja edycji} (\code{validateDraft}): oba terminy muszą być ustawione, koniec musi być po
starcie, a start w przyszłości. Typy \code{RoomEventState}, \code{RsvpEntry}, \code{RsvpStatus} pochodzą z
\code{chatStore}; akcje \code{?/scheduleEvent}, \code{?/cancelEvent}, \code{?/setRsvp} opisano w rozdziale
o trasach.

\section{Komponenty regionalne}\label{sec:fe-komp-regiony}

\begin{longtable}{@{}p{0.24\textwidth}p{0.30\textwidth}p{0.38\textwidth}@{}}
\toprule
\textbf{Komponent} & \textbf{Propsy} & \textbf{Rola} \\
\midrule
\endhead
\code{RegionChooser} & \code{regions: RegionConfig[]} & Ekran wyboru regionu na hoście wejściowym (pomiar, failover, Stay/Switch) \\
\code{RegionProposalBanner} & \code{currentRegion: RegionId}, \code{other: RegionConfig | undefined} & Baner „Keep current / Switch” na hoście regionalnym \\
\bottomrule
\end{longtable}

Pełny opis logiki znajduje się w sekcji „Regiony w interfejsie” rozdziału o bibliotece klienta.
Oba komponenty składają się z elementów UI kitu: \code{Crest}, \code{InnerPanel}, \code{OrnateButton},
\code{SectionTitle}.

\section{UI kit w stylu Diablo II (\texttt{chrome})}\label{sec:fe-chrome}

Zestaw 18 widżetów w \code{frontend/src/lib/components/chrome/}. Komponenty przyjmują
\code{children?: import('svelte').Snippet} tam, gdzie wskazano, i używają \code{$props()} z
\code{interface Props}. Poniższa tabela zbiera propsy; kolumna „B” oznacza prop bindowalny
(\code{$bindable}).

\small
\begin{longtable}{@{}p{0.15\textwidth}p{0.29\textwidth}p{0.11\textwidth}p{0.03\textwidth}p{0.32\textwidth}@{}}
\toprule
\textbf{Prop} & \textbf{Typ} & \textbf{Domyśl.} & \textbf{B} & \textbf{Opis} \\
\midrule
\endhead
\multicolumn{5}{@{}l}{\textbf{Cycler} (generyczny \code{<T>}) --- przełączanie wartości z tablicy strzałkami, z zawijaniem na krańcach} \\
\code{values} & \code{T[]} & --- & & Pula elementów \\
\code{value} & \code{T} & --- & tak & Wybrana wartość \\
\code{label} & \code{string | undefined} & --- & & Opcjonalna etykieta \\
\code{format} & \code{(v: T) => string} & --- & & Własny formater wyświetlania \\
\code{disabled} & \code{boolean} & \code{false} & & Wyłącza nawigację \\
\code{onchange} & \code{(v: T) => void} & --- & & Wywołanie z nową wartością \\
\midrule
\multicolumn{5}{@{}l}{\textbf{Crest} --- tarcza SVG w stylu D2 (celtyckie sploty, gwiazda kompasu, nity)} \\
\code{size} & \code{number} & \code{96} & & Szerokość/wysokość w px \\
\code{tone} & \code{'gold' | 'bone' | 'iron'} & \code{'gold'} & & Wariant kolorystyczny \\
\midrule
\multicolumn{5}{@{}l}{\textbf{D2DatePicker} --- kalendarz w popoverze (siatka 6$\times$7 = 42 komórki)} \\
\code{value} & \code{Date | null} & --- & & Wybrana data \\
\code{onChange} & \code{(d: Date) => void} & --- & & Po wyborze daty (zachowuje porę dnia) \\
\code{min} & \code{Date | null} & \code{null} & & Minimalna wybieralna data \\
\code{placeholder} & \code{string} & \code{'Pick a date'} & & Tekst bez wybranej daty \\
\code{ariaLabel} & \code{string} & \code{'Date'} & & \code{aria-label} przycisku \\
\midrule
\multicolumn{5}{@{}l}{\textbf{D2TimePicker} --- popover z dwiema przewijanymi kolumnami (godziny 00--23, minuty)} \\
\code{value} & \code{string} & --- & & Czas \code{"HH:MM"} (24 h) \\
\code{onChange} & \code{(value: string) => void} & --- & & Wywołanie z \code{"HH:MM"} \\
\code{step} & \code{5 | 10 | 15 | 30} & \code{5} & & Granulacja minut \\
\code{ariaLabel} & \code{string} & \code{'Time'} & & \code{aria-label} przycisku \\
\midrule
\multicolumn{5}{@{}l}{\textbf{Hint} --- plakietka podpowiedzi klawisza} \\
\code{key} & \code{string} & --- & & Znak klawisza (np. \code{"E"}, \code{"M"}) \\
\code{label} & \code{string} & --- & & Opis \\
\code{tone} & \code{'gold' | 'green' | 'amber' | 'red' | 'stone' | 'blue'} & \code{'stone'} & & Akcent kolorystyczny \\
\code{onclick} & \code{() => void} & --- & & Ustawiony: renderuje \code{<button>}, inaczej \code{<span>} \\
\midrule
\multicolumn{5}{@{}l}{\textbf{HintBar} --- stopka z podpowiedziami (ukryta na urządzeniach dotykowych)} \\
\code{children} & \code{Snippet} & --- & & Elementy podpowiedzi (zwykle \code{Hint}) \\
\midrule
\multicolumn{5}{@{}l}{\textbf{InnerPanel} --- kamienna płyta z brązową ramą, czterema nitami i szumem ziarna} \\
\code{title} & \code{string | undefined} & --- & & Nagłówek panelu \\
\code{padded} & \code{boolean} & \code{true} & & Wewnętrzny margines \\
\code{children} & \code{Snippet} & --- & & Treść panelu \\
\code{actions} & \code{Snippet} & --- & & Akcje obok tytułu \\
\midrule
\multicolumn{5}{@{}l}{\textbf{ListRow} --- wiersz siatki list pokoi i członków} \\
\code{selected} & \code{boolean} & \code{false} & & Zaznaczony (złoty pip) \\
\code{header} & \code{boolean} & \code{false} & & Wiersz nagłówka (bez interakcji) \\
\code{disabled} & \code{boolean} & \code{false} & & Wyłączony \\
\code{member} & \code{boolean} & \code{false} & & Stan członka (przygaszony pip) \\
\code{onclick} & \code{() => void} & --- & & Czyni wiersz interaktywnym \\
\code{children} & \code{Snippet} & --- & & Treść \\
\midrule
\multicolumn{5}{@{}l}{\textbf{OrnateButton} --- przycisk lub link w wielu wariantach} \\
\code{variant} & \code{'primary' | 'secondary' | 'danger' | 'ghost'} & \code{'secondary'} & & Styl \\
\code{size} & \code{'sm' | 'md' | 'lg'} & \code{'md'} & & Rozmiar \\
\code{type} & \code{'button' | 'submit' | 'reset'} & \code{'button'} & & Ignorowany przy \code{href} \\
\code{disabled} & \code{boolean} & \code{false} & & Wyłączony \\
\code{loading} & \code{boolean} & \code{false} & & Ładowanie (też wyłącza) \\
\code{href} & \code{string | undefined} & --- & & Ustawiony: \code{<a role="button">} \\
\code{fullWidth} & \code{boolean} & \code{false} & & Rozciąga na szerokość kontenera \\
\code{onclick} & \code{(e: MouseEvent) => void} & --- & & Obsługa kliknięcia \\
\code{children} & \code{Snippet} & --- & & Etykieta \\
\midrule
\multicolumn{5}{@{}l}{\textbf{PageFrame} --- ramka strony: kamienny mur z kwadrolistem, winieta, cztery rogi SVG, złote krawędzie} \\
\code{children} & \code{Snippet} & --- & & Treść strony (wewnątrz \code{.frame-inner}) \\
\midrule
\multicolumn{5}{@{}l}{\textbf{PipMeter} --- rząd orbów pokazujący liczność} \\
\code{value} & \code{number} & --- & & Wartość (zaokrąglana w dół do liczby zapalonych pipsów) \\
\code{total} & \code{number} & \code{3} & & Łączna liczba pipsów \\
\code{tone} & \code{'good' | 'mid' | 'bad' | 'auto'} & \code{'auto'} & & Schemat kolorów \\
\code{size} & \code{'sm' | 'md'} & \code{'md'} & & Rozmiar pipsa \\
\code{glow} & \code{boolean} & \code{true} & & Poświata zapalonych \\
\code{label} & \code{string | undefined} & --- & & Etykieta pod pipsami \\
\midrule
\multicolumn{5}{@{}l}{\textbf{ProgressBar} --- pasek postępu z opcjonalnymi znacznikami} \\
\code{value} & \code{number} & --- & & Wartość bieżąca \\
\code{max} & \code{number} & \code{100} & & Maksimum \\
\code{height} & \code{number} & \code{6} & & Wysokość wypełnienia (px) \\
\code{ticks} & \code{boolean} & \code{false} & & Znaczniki 25/50/75\% \\
\code{variant} & \code{'gold' | 'blood' | 'green'} & \code{'gold'} & & Kolor wypełnienia \\
\code{label} & \code{string | undefined} & --- & & Etykieta i licznik \\
\midrule
\multicolumn{5}{@{}l}{\textbf{SectionTitle} --- nagłówek sekcji ze złotą linią i slotem z prawej} \\
\code{title} & \code{string} & --- & & Tekst nagłówka \\
\code{size} & \code{'small' | 'normal' | 'large'} & \code{'normal'} & & Stopień wielkości \\
\code{tone} & \code{'gold' | 'bone' | 'blood'} & \code{'gold'} & & Kolor tytułu \\
\code{children} & \code{Snippet} & --- & & Zawartość slotu za linią \\
\midrule
\multicolumn{5}{@{}l}{\textbf{Sigil} --- deterministyczny sigil SVG z adresu Ethereum} \\
\code{address} & \code{string} & --- & & Adres z prefiksem \code{0x} \\
\code{size} & \code{number} & \code{56} & & Szerokość/wysokość w px \\
\midrule
\multicolumn{5}{@{}l}{\textbf{StoneCheckbox} --- pole wyboru ze „kamiennym” polem i animowanym znaczkiem} \\
\code{checked} & \code{boolean} & \code{false} & tak & Stan zaznaczenia \\
\code{label} & \code{string | undefined} & --- & & Etykieta (kapitaliki) \\
\code{disabled} & \code{boolean} & \code{false} & & Wyłączone \\
\code{size} & \code{'sm' | 'md'} & \code{'md'} & & Rozmiar \\
\code{onchange} & \code{(checked: boolean) => void} & --- & & Wywołanie przy zmianie \\
\midrule
\multicolumn{5}{@{}l}{\textbf{SystemDialog} --- okno modalne, wbudowane lub pływające} \\
\code{open} & \code{boolean} & \code{false} & tak & Widoczność \\
\code{title} & \code{string} & --- & & Tytuł \\
\code{tone} & \code{'gold' | 'blood' | 'green' | 'iron'} & \code{'iron'} & & Akcent kolorystyczny \\
\code{modal} & \code{boolean} & \code{false} & & Ciemna nakładka, Escape zamyka, blokada przewijania \\
\code{inline} & \code{boolean} & \code{false} & & Render w przepływie dokumentu \\
\code{position} & \code{'bottom-right' | 'top-center' | 'center'} & \code{'bottom-right'} & & Położenie okna pływającego \\
\code{width} & \code{string} & \code{'420px'} & & Szerokość (CSS) \\
\code{closeable} & \code{boolean} & \code{true} & & \code{false} ukrywa przycisk zamknięcia i blokuje zamykanie \\
\code{onclose} & \code{() => void} & --- & & Wywołanie przy zamknięciu \\
\code{children} & \code{Snippet} & --- & & Treść \\
\code{footer} & \code{Snippet} & --- & & Stopka pod treścią \\
\midrule
\multicolumn{5}{@{}l}{\textbf{Tab} --- zakładka górnego paska nawigacji (aktywność z \code{$app/state})} \\
\code{href} & \code{string} & --- & & Ścieżka trasy \\
\code{match} & \code{'exact' | 'prefix'} & \code{'prefix'} & & Strategia dopasowania \\
\code{label} & \code{string | undefined} & --- & & Tekst zakładki \\
\code{children} & \code{Snippet} & --- & & Alternatywa dla \code{label} \\
\midrule
\multicolumn{5}{@{}l}{\textbf{Tabs} --- kontener paska zakładek (\code{grid-auto-flow: column}, złote linie u góry i u dołu)} \\
\code{children} & \code{Snippet} & --- & & Elementy \code{Tab} \\
\bottomrule
\end{longtable}
\normalsize

\subsection{Zachowanie widżetów}

\begin{itemize}[nosep]
  \item \textbf{Cycler:} \code{$effect} inicjuje \code{value} wartością \code{values[0]}, jeśli \code{undefined};
        strzałki pulsują 120 ms po kliknięciu.
  \item \textbf{Crest:} renderowany jako \code{<span aria-hidden="true" role="presentation">} z wbudowanym SVG.
  \item \textbf{D2DatePicker:} tydzień zaczyna się od poniedziałku; popover odwraca się nad pole, gdy pod nim
        jest mniej niż 320 px; kliknięcie komórki spoza bieżącego miesiąca przechodzi do tego miesiąca;
        przycisk „Today” (\code{jumpToday()}); zamyka się po kliknięciu poza oknem i klawiszem Escape; cztery
        brązowe nity w rogach (zgodnie z estetyką \code{InnerPanel}).
  \item \textbf{D2TimePicker:} parsuje \code{HH:MM}, a przy błędzie przyjmuje \code{18:00}; wybrana godzina i
        minuta przewijają się do widoku po otwarciu; zamyka się po kliknięciu poza oknem i Escape.
  \item \textbf{Hint / HintBar:} typ \code{HintTone} jest wspólny z \code{hintsContext}; większość tras
        przekazuje \code{HintEntry[]} przez \code{provideHints()} / \code{getHintsState()}.
        \code{HintBar} renderuje się jako \code{<footer aria-label="Keyboard hints">} i jest ukryty regułą
        \code{@media (hover: none) and (pointer: coarse)}.
  \item \textbf{InnerPanel:} gdy podano \code{title} lub \code{actions}, nad treścią powstaje nagłówek ze
        złotą linią; ziarno kamienia to pseudoelement CSS z warstwowym szumem.
  \item \textbf{ListRow:} gdy \code{!header && !disabled && typeof onclick === 'function'}, wiersz dostaje
        \code{role="button"}, \code{tabindex="0"} i reaguje na Enter/Spację; kliknięcie wywołuje animację
        pulsu 180 ms; \code{aria-selected} i \code{aria-disabled} ustawiane są automatycznie.
  \item \textbf{OrnateButton:} przy \code{loading} wewnętrzny \code{ornate__inner} pulsuje animacją CSS;
        \code{fullWidth} dodaje klasę z \code{width: 100\%}.
  \item \textbf{PageFrame:} renderowany jako \code{<div class="page-frame anim-frame-in">}; warstwa tła to
        ciepły ciemnobrązowy mur z powtarzalnym wzorem kwadrolistu i radialną winietą; cztery rogi SVG
        w warstwie chrome używają złotego gradientu obrysu, a cztery \code{<span>} tworzą włosowe krawędzie.
  \item \textbf{PipMeter, ton automatyczny:} 3 zapalone lub więcej $\rightarrow$ \code{good} (zielony); 2
        $\rightarrow$ \code{mid} (bursztynowy); 1 $\rightarrow$ \code{bad} (czerwony); 0 $\rightarrow$
        \code{good} (domyślnie). Zapalające się pipsy mają animację \code{orbPop} 300 ms.
  \item \textbf{ProgressBar:} szerokość wypełnienia to \code{clamped / safeMax * 100\%} z
        \code{transition: width 0.35s ease}; \code{safeMax} zamienia \code{max <= 0} na 1; ustawiane są rola
        \code{progressbar} oraz \code{aria-valuenow}/\code{aria-valuemax}.
  \item \textbf{SectionTitle:} renderuje \code{<h3 class="title small-caps">}; linia to elastyczna włosowa kreska
        o grubości 1 px.
  \item \textbf{Sigil:} algorytm \code{parseSpec}: (1) adres małymi literami bez \code{0x}; (2) wymagane co
        najmniej 12 znaków hex, uzupełnienie deterministyczne do 16 bajtów; (3) \code{bytes[0] \% 6} to indeks
        kształtu (0--5); (4) \code{bytes[1] \% 3} to typ symetrii: 0 --- odbicie pionowe, 1 --- obrót
        4-krotny, 2 --- obrót 3-krotny; (5) \code{bytes[10] \% 2 === 1} włącza złoty akcent; (6)
        \code{bytes[11] & 0x0f} wybiera pierścień akcentu; (7) 25 komórek (siatka 5$\times$5) z
        \code{(bytes[2 + i*4] >> (i \% 8)) & 1}, następnie składanych według symetrii. Wynik to
        \code{<div class="sigil-frame bevel-out" aria-hidden="true">} z SVG; nieprawidłowy adres daje awaryjną
        ikonę skrzyżowanych mieczy (wzór \code{X}).
  \item \textbf{StoneCheckbox:} dla dostępności używany jest prawdziwy \code{<input type="checkbox" class="sr-only">};
        wizualne pole to \code{<span>} z animowanym SVG; całość opakowuje \code{<label>}.
  \item \textbf{SystemDialog:} cztery tryby: modalny (nakładka, centrowanie, \code{overflow: hidden} na body,
        Escape zamyka), wbudowany (\code{inline}), pływający (domyślny, pozycjonowany według \code{position}),
        ukryty (\code{open: false}, nic nie renderuje). Ton steruje kolorem linii nagłówka, przycisku zamknięcia
        i odcieniem nakładki; zamknięcie ustawia \code{open = false} i wywołuje \code{onclose}.
  \item \textbf{Tab:} dla \code{match = 'prefix'} jest aktywna, gdy \code{page.url.pathname === href} lub zaczyna się
        od \code{href + '/'}; dla \code{'exact'} liczy się tylko równość. Aktywna zakładka ma złotą strzałkę pod
        tekstem; używa \code{data-sveltekit-preload-data="hover"}.
  \item \textbf{Tabs:} przy szerokości widoku $\le$ 600 px wiersz ma \code{overflow-x: auto} z ukrytym paskiem
        przewijania; górna i dolna linia to gradienty przezroczysty $\rightarrow$ złoty $\rightarrow$ przezroczysty.
\end{itemize}

\section{Odnośniki między warstwami}\label{sec:fe-odnosniki}

\begin{itemize}[nosep]
  \item Endpointy REST wołane z form actions oraz ramki WebSocket (\code{/ws/rooms}, czat pokoju) opisują
        rozdziały o backendzie (REST API i protokoły realtime).
  \item Encje \code{Room}, \code{RoomEvent}, \code{RoomMember} opisuje rozdział o modelu danych.
  \item Wdrożenie wieloregionowe, zmienne \code{PUBLIC_REGION*}, sondy \code{/healthz} i log VictoriaLogs ---
        część o operacjach i architekturze.
\end{itemize}

% SPRZECZNOSC: docs/frontend/routes.md i frontend/README.md opisują ciasteczko sesji jako
% sameSite=strict; kod (routes/auth/+page.server.ts) ustawia sameSite=lax (komentarz: linki z Discorda).
% Przyjęto kod (nowsze). Pozostałe różnice względem docs: ping na /rooms woła /healthz (docs: GET /);
% trasa /test nadal woła GET /; loadery/akcje używają backendFetch (docs: bezpośredni $BACKEND_BASE);
% README.md nie wymienia tras /about, /profile, /rooms/[name] ani regionów.
